%% notes-tex/010-additive-baselines/sections/3-what-each-model-sees.tex

\section{What each model actually sees\secstatus{todo}}
\label{sec:sees}

The comparison is only meaningful if the models are given the same information.
Reading the 010 code settles that, and the answer is stronger than expected:
every model in this document, the transformer included, receives exactly one
thing per record, the set of three perturbed genes.

\subsection{The transformer's per-sample input\secstatus{todo}}

The 010 dataset is built with the \texttt{Perturbation} graph processor. That
processor writes the perturbed gene names, their indices, a boolean mask and the
label into each sample, and nothing else
(\texttt{graph\_processor.py}, lines 1984 to 2058). There is no per-gene feature
matrix: with an empty \texttt{node\_embeddings} list in the config, the gene
feature tensor of the cell graph has shape $6607 \times 0$
(\texttt{cell\_data.py}, line 53). Gene identity is carried by a free lookup
table of 6{,}607 rows and 180 columns
(\texttt{equivariant\_cell\_graph\_transformer.py}, line 1989), initialized
randomly and fit by gradient descent.

No sequence, expression, ontology or other precomputed representation enters.
The model is given gene identities, the same as a one-hot design matrix, and
learns its own 180-dimensional vector for each.

\subsection{The encoder does not depend on the strain\secstatus{todo}}
\label{sec:encoder}

The forward pass builds its input as
\texttt{gene\_embedding(torch.arange(N))}, prepends a CLS token, and unsqueezes
to shape $[1, N{+}1, 180]$. That leading dimension is $1$, not the batch size.
The eight transformer layers then run on this single tensor, which is assembled
entirely from parameters and contains nothing from \texttt{batch}. Writing $E$
for the embedding table and $T$ for the eight-layer encoder,
%
\begin{equation}
  Z \;=\; T(E) \in \mathbb{R}^{N \times 180},
  \qquad
  h_{\mathrm{CLS}} \in \mathbb{R}^{180},
  \label{eq:encoder}
\end{equation}
%
and $Z$ and $h_{\mathrm{CLS}}$ are the same for every record in the dataset. The
repository states this itself, in the docstring of
\texttt{PerturbationGraphPropagation}: ``The encoder runs at batch 1 on the
WILDTYPE graph, so \texttt{H\_genes} and \texttt{h\_CLS} are identical for every
strain; the only strain-dependent step is
\texttt{EquivariantPerturbationTransform}.''

So the eight encoder layers, which hold 3{,}129{,}120 of the model's 4{,}774{,}861
parameters, perform no per-sample computation. They map one free parameter block
to another. Whatever role they play, it is not to process the strain.

\subsection{Where the nine graphs act, and where they do not\secstatus{todo}}
\label{sec:graphs}

The nine adjacency matrices enter the model in exactly one place. During
training, a Kullback-Leibler penalty pulls one attention head of layer 1 toward
each row-normalized adjacency $\tilde{A}^{(k)}$,
%
\begin{equation}
  \mathcal{L}_{\mathrm{graph}} \;=\; \sum_{k=1}^{9} c_k
  \sum_{i=1}^{N} \sum_{j=1}^{N}
  \tilde{A}^{(k)}_{ij} \log \frac{\tilde{A}^{(k)}_{ij}}{\alpha^{(k)}_{ij}},
  \qquad
  \alpha^{(k)} = \mathrm{softmax}\!\left(\frac{QK^{\top}}{\sqrt{d}}\right)_{\!h_k},
  \label{eq:graphreg}
\end{equation}
%
and that is the only operation that reads them.

\notesec{Not in the prediction}
Three mechanisms could put adjacency into the computation that yields
$\hat{y}$: a head mask on attention, the perturbation-propagation block, and an
adjacency bias inside the attention logits. All three are off in the 010
configuration. Given a trained checkpoint, no adjacency tensor is an input to any
operation between the gene indices and the prediction, and
Section~\ref{sec:verify} checks this directly: setting the penalty weight to zero
on a trained checkpoint moves the predictions by at most $1.3 \times 10^{-6}$.

\notesec{In the weights}
The penalty is computed inside the same forward call as the prediction, from the
same pre-dropout attention tensor, and nothing detaches it. Its gradient reaches
the layer-1 query and key projections, through them every layer-0 parameter,
and through those the gene embedding table. So $Z = T(E)$ of
Eq.~\ref{eq:encoder}, the per-gene vector the prediction is read from, is shaped
by the nine graphs during training even though its value at inference comes off
fixed weights. The graphs act on the parameters, not on the function from
checkpoint to prediction. An ablation that zeroes the coefficient on a trained
checkpoint tests only the second.

\notesec{How large the penalty was}
The per-graph weight $c_k$ in Eq.~\ref{eq:graphreg} is a product of three
factors: the per-head $\lambda_k = 10^{-3}$, a global scale of $1.0$ read from
\texttt{regression\_task.loss.graph\_regularization.lambda}, and an edge
normalization meant to divide by the mean degree. That last factor is computed as
\texttt{len(A.nonzero()[0])} summed over the nine matrices. On a dense
\texttt{torch} tensor \texttt{A.nonzero()} has shape $[\mathrm{nnz}, 2]$, so
\texttt{[0]} is one index pair of length two, and the sum over nine graphs is
$18$ rather than the 2.5 million edges intended. Dividing by $18/6607$ multiplies
by $367.06$,
%
\begin{equation}
  c_k \;=\; 10^{-3} \times 1.0 \times \frac{6607}{18} \;=\; 0.367 ,
  \label{eq:effcoef}
\end{equation}
%
applied to a divergence summed over all $6{,}607$ rows. The runs' logs record
the consequence: against a point loss of $0.81$ the graph term is $5{,}740$, and
the trainer's diagnostic \texttt{norm\_weighted\_graph\_reg}, the graph term
divided by the total loss, reads $0.99986$ for all three checkpoints. The
optimizer spent 99.99 percent of its objective on the graph penalty.
Figure~\ref{fig:penalty} shows the three loss terms and that share over the
whole of training. Section~\ref{sec:graphablation} reports what the penalty
bought: the regularized heads recover their graphs, and the matched runs without
the penalty do not train.

The dominance was not what produced the accuracy. The 025 replication of this
configuration (job 1598, Section~\ref{sec:disjoint025}) ran after the edge
normalization had been corrected, so its graph term is $0.34$ against a point
loss of $0.80$, a share of 0.29 to 0.33 rather than 0.9999, and it reached
validation Pearson 0.4463 at epoch 14 against 0.4472 to 0.4619 for the three
010 checkpoints.

%% SOURCE: experiments/010-kuzmin-tmi/scripts/graph_penalty_vs_loss.py ->
%% results/graph_penalty_vs_loss_history.csv, pulled from the four runs' W&B
%% histories; summary in results/graph_penalty_vs_loss_summary.json
\tcfigfit[0.98]{figures/graph_penalty_vs_loss.pdf}{The training objective of
the three 010 runs and of the 025 replication, per epoch. \textbf{a}, the point
term, the Sinkhorn term at its weight of 0.1, and the weighted graph penalty for
the 010 runs, on a log axis: the penalty sits four orders of magnitude above the
point loss from epoch 5 on. \textbf{b}, the same three terms for job 1598, where
the corrected edge normalization puts the penalty below the point loss.
\textbf{c}, the graph term's share of the total loss, recomputed from the logged
terms: 0.9999 throughout for the 010 runs, 0.32 for the replication.
\textbf{d}, validation Pearson per epoch with the best-Pearson epoch marked; the
replication reaches the 010 band with the penalty at a third of the
objective.}{fig:penalty}

All nine heads carried the penalty. The configuration names the first two graphs
\texttt{physical} and \texttt{regulatory}, which is what the cell graph called
them at the December 2025 commit these runs used. The suffixed names,
\texttt{physical\_interaction} and \texttt{regulatory\_interaction}, date from
2026-06-27. The runs log the edge-recovery keys \texttt{physical\_L1\_H0} and
\texttt{regulatory\_L1\_H1}, which the trainer emits only for a graph it
matched.
An earlier version of this document claimed those two heads were unregularized
and that the weight was squared to $10^{-6}$; both claims were checked against
the run logs and are false for these runs.

\subsection{The transformer is a set function on three genes\secstatus{todo}}
\label{sec:setfunction}

What remains after Eq.~\ref{eq:encoder} is the strain-dependent part. For a
record with perturbed set $S$, $|S| = 3$, the perturbation transform performs
cross-attention in which all $N$ genes query the $|S|$ rows $\{Z_p : p \in S\}$,
followed by a residual, a layer norm and a feedforward block, giving
$H^{\mathrm{pert}}_i = g(Z_i + c_S)$ where $c_S$ is the attended context. The
head then pools over the perturbed genes and reads out,
%
\begin{equation}
  z_S \;=\; \frac{1}{|S|}\sum_{i \in S} H^{\mathrm{pert}}_i ,
  \qquad
  \hat{y} \;=\; \mathrm{MLP}\!\left([\,h_{\mathrm{CLS}} \,\Vert\, z_S\,]\right),
  \label{eq:cgt}
\end{equation}
%
with the MLP being $\mathrm{Linear}(360, 180) \to \mathrm{ReLU} \to
\mathrm{Linear}(180, 1)$.

Since $h_{\mathrm{CLS}}$ is constant, Eq.~\ref{eq:cgt} is a function of $z_S$
alone, and $z_S$ is a permutation-invariant pooling of quantities derived from
$\{Z_p : p \in S\}$. The 010 transformer is therefore a set function on the
three perturbed gene identities, built on a free per-gene vector. That is the
same functional shape as the additive nulls, which are set functions on the same
three identities built on a free per-gene scalar. The models differ in how
expressive the set function is, not in what they are allowed to see.

This is what makes the comparison in Section~\ref{sec:results} fair. It does not
settle where the transformer's advantage over the nonlinear baseline comes from.
An earlier version of the working note assigned that gap to the nine interaction
graphs, and a later one assigned it to capacity on the grounds that the graphs
are absent from the prediction. Both overreached. The graphs shaped the weights
this set function is built from, and Section~\ref{sec:graphablation} measures how
heavily.
